\import{set.tex}
\import{set/cons.tex}
\import{set/powerset.tex}
\import{set/regularity.tex}
\import{set/suc.tex}


\section{Transitive sets}

We use the word \emph{transitive} to talk about sets as relations,
so we will explicitly talk about \emph{\in-transitivity} here.

The definition of \in-transitivity is supplied directly by the standard
prelude.

\begin{proposition}\label{transitiveset_iff_subseteq}
    $A$ is \in-transitive iff
    for all $a\in A$ we have $a\subseteq A$.
\end{proposition}
\begin{proof}
    Follows by \auto{transitiveset,subseteq}.
\end{proof}

\begin{proposition}\label{transitiveset_iff_pow}
    $A$ is \in-transitive iff $A\subseteq \pow{A}$.
\end{proposition}
\begin{proof}
    For all $a\in A$ we have $a\subseteq A \iff a\in\pow{A}$
        by \obv{subseteq,pow_iff}.
    Follows by \auto{elem_subseteq,subseteq,pow_iff,transitiveset_iff_subseteq}.
\end{proof}

\begin{proposition}\label{transitiveset_iff_unions_suc}
    $A$ is \in-transitive iff $\unions{\suc{A}} = A$.
\end{proposition}
\begin{proof}
    We have if $A$ is \in-transitive, then $\unions{\suc{A}}=A$
        by \auto{setext,transitiveset,unions_iff,suc_iff}.
    We have if $\unions{\suc{A}}=A$, then $A$ is \in-transitive
        by \auto{transitiveset,unions_iff,suc_iff}.
    Follows by \auto.
\end{proof}

\begin{proposition}\label{transitiveset_iff_unions_subseteq}
    $A$ is \in-transitive iff $\unions{A}\subseteq A$.
\end{proposition}
\begin{proof}
    Follows by \auto{transitiveset,subseteq,unions_iff}.
\end{proof}

\begin{proposition}\label{transitiveset_upair}
    Suppose $A$ is \in-transitive and $\{a,b\}\in A$.
    Then $a,b\in A$.
\end{proposition}
\begin{proof}
    Follows by \auto{transitiveset,transitiveset_iff_subseteq,subseteq}.
\end{proof}

\subsection{Closure properties of \in-transitive sets}

\begin{proposition}\label{emptyset_transitiveset}
    $\emptyset$ is \in-transitive.
\end{proposition}
\begin{proof}
    Follows by \obv{transitiveset,emptyset}.
\end{proof}

\begin{proposition}\label{union_of_transitiveset_is_transitiveset}
    Suppose $A$ and $B$ are \in-transitive.
    Then $A\union B$ is \in-transitive.
\end{proposition}
\begin{proof}
    Follows by \auto{transitiveset,union_iff}.
\end{proof}

\begin{proposition}\label{inter_of_transitiveset_is_transitiveset}
    Let $A, B$ be \in-transitive sets.
    Then $A\inter B$ is \in-transitive.
\end{proposition}
\begin{proof}
    Follows by \obv{transitiveset,inter}.
\end{proof}

\begin{proposition}\label{suc_of_transitiveset_is_transitiveset}
    Let $A$ be an \in-transitive set.
    Then $\suc{A}$ is \in-transitive.
\end{proposition}
\begin{proof}
    Follows by \auto{transitiveset,suc}.
\end{proof}

\begin{proposition}\label{unions_of_transitiveset_is_transitiveset}
    Let $A$ be an \in-transitive set.
    Then $\unions{A}$ is \in-transitive.
\end{proposition}
\begin{proof}
    Follows by \auto{transitiveset,unions_iff}.
\end{proof}

\begin{proposition}\label{unions_family_of_transitiveset_is_transitiveset}
    Suppose every element of $A$ is an \in-transitive set.
    Then $\unions{A}$ is \in-transitive.
\end{proposition}
\begin{proof}
    Follows by \auto{transitiveset,unions_iff}.
\end{proof}

\begin{proposition}\label{inters_family_of_transitiveset_is_transitiveset}
    Suppose every element of $A$ is an \in-transitive set.
    Then $\inters{A}$ is \in-transitive.
\end{proposition}
\begin{proof}
    Follows by \auto{transitiveset,inters,unions_family_of_transitiveset_is_transitiveset}.
\end{proof}


\section{Ordinals}

\begin{definition}\label{ordinal}
    $\alpha$ is an ordinal iff
    $\alpha$ is \in-transitive and
    every element of $\alpha$ is \in-transitive.
\end{definition}

\begin{proposition}\label{ordinal_intro}
    Suppose $\alpha$ is \in-transitive and
        every element of $\alpha$ is \in-transitive.
    Then $\alpha$ is an ordinal.
\end{proposition}
\begin{proof}
    Follows by \obv{ordinal}.
\end{proof}

\begin{proposition}\label{ordinal_is_transitiveset}
    Let $\alpha$ be an ordinal.
    Then $\alpha$ is \in-transitive.
\end{proposition}
\begin{proof}
    Follows by \obv{ordinal}.
\end{proof}

\begin{proposition}\label{ordinal_elem_is_transitiveset}
    Let $\alpha$ be an ordinal.
    Suppose $A\in\alpha$.
    Then $A$ is \in-transitive.
\end{proposition}
\begin{proof}
    Follows by \obv{ordinal}.
\end{proof}

\begin{classification}\label{elem_of_ordinal_is_ordinal}
    Let $\alpha$ be an ordinal.
    Suppose $\beta\in\alpha$.
    Then $\beta$ is an ordinal.
\end{classification}
\begin{proof}
    Follows by \auto{ordinal,transitiveset}.
\end{proof}

\begin{proposition}\label{suc_ordinal_implies_ordinal}
    Suppose $\suc{\alpha}$ is an ordinal.
    Then $\alpha$ is an ordinal.
\end{proposition}
\begin{proof}
    Follows by \auto{ordinal,suc}.
\end{proof}

\begin{proposition}\label{transitivesubseteq_of_ordinal_is_ordinal}
    Let $\alpha$ be an ordinal.
    Suppose $\beta\subseteq\alpha$.
    Suppose $\beta$ is \in-transitive.
    Then $\beta$ is an ordinal.
\end{proposition}
\begin{proof}
    Follows by \obv{subseteq,ordinal}.
\end{proof}

\begin{proposition}\label{ordinal_elem_implies_subseteq}
    Let $\alpha,\beta$ be ordinals.
    Suppose $\alpha\in\beta$.
    Then $\alpha\subseteq\beta$.
\end{proposition}
\begin{proof}
    Follows by \auto{ordinal,transitiveset_iff_subseteq}.
\end{proof}

\begin{proposition}\label{ordinal_transitivity}
    Let $\alpha$ be an ordinal.
    Suppose $\gamma\in\beta\in\alpha$.
    Then $\gamma\in\alpha$.
\end{proposition}
\begin{proof}
    Follows by \auto{ordinal,transitiveset}.
\end{proof}

\begin{proposition}\label{ordinal_suc_subseteq}
    Let $\beta$ be an ordinal.
    Suppose $\alpha\in\beta$.
    Then $\suc{\alpha}\subseteq\beta$.
\end{proposition}
\begin{proof}
    Follows by \auto{ordinal,transitiveset_iff_subseteq,suc_subseteq_intro}.
\end{proof}


\begin{abbreviation}\label{ordinal_prec}
    $\alpha \precedes \beta$ iff
    $\beta$ is an ordinal and
    $\alpha\in\beta$.
\end{abbreviation}

\begin{abbreviation}\label{ordinal_preceq}
    $\alpha\precedeseq \beta$ iff $\beta$ is an ordinal and $\alpha\subseteq\beta$.
\end{abbreviation}

\begin{lemma}\label{prec_is_ordinal}
    Let $\alpha,\beta$ be sets.
    Suppose $\alpha\precedes \beta$.
    Then $\alpha$ is an ordinal.
\end{lemma}
\begin{proof}
    Follows by \obv{elem_of_ordinal_is_ordinal}.
\end{proof}

We already have global irreflexivity and asymmetry of \in.
\in\ is transitive on ordinals by definition.
To show that \in\ is a strict total order it only remains to show that \in\ is connex.

\begin{lemma}\label{ordinal_elem_connex_from_elements}
    Let $\alpha,\gamma$ be ordinals.
    Suppose for all sets $\xi,\beta$ such that $\xi\in\alpha$ and $\beta$ is an ordinal
    we have $\xi\in\beta\lor \beta\in\xi \lor \xi = \beta$.
    Suppose for every $\eta\in\gamma$
    we have $\alpha\in\eta\lor \eta\in\alpha \lor \alpha = \eta$.
    Then $\alpha\in\gamma\lor \gamma\in\alpha \lor \alpha = \gamma$.
\end{lemma}
\begin{proof}
    Assume neither $\alpha\in\gamma$ nor $\gamma\in\alpha$.
    For every $\xi\in\alpha$ we have $\xi\in\gamma$
        by \obv{ordinal_transitivity}.
    For every $\eta\in\gamma$ we have $\eta\in\alpha$
        by \obv{ordinal_transitivity}.
    Follows by set extensionality.
\end{proof}

\begin{lemma}\label{ordinal_elem_connex_step}
    Let $\alpha,\gamma$ be ordinals.
    Suppose for all sets $\xi,\beta$ such that $\xi\in\alpha$ and $\beta$ is an ordinal
    we have $\xi\in\beta\lor \beta\in\xi \lor \xi = \beta$.
    Then $\alpha\in\gamma\lor \gamma\in\alpha \lor \alpha = \gamma$.
\end{lemma}
\begin{proof}[Proof by \in-induction on $\gamma$]
    Follows by \obv{ordinal_elem_connex_from_elements,elem_of_ordinal_is_ordinal}.
\end{proof}

\begin{proposition}\label{ordinal_elem_connex}
    Let $\alpha$ be an ordinal.
    Then for every ordinal $\beta$
    we have $\alpha\in\beta\lor \beta\in\alpha \lor \alpha = \beta$.
\end{proposition}
\begin{proof}[Proof by \in-induction on $\alpha$]
    Follows by \obv{ordinal_elem_connex_step,elem_of_ordinal_is_ordinal}.
\end{proof}

\begin{proposition}\label{ordinal_proper_subset_implies_elem}
    Let $\alpha,\beta$ be ordinals.
    Suppose $\alpha\subset\beta$.
    Then $\alpha\in\beta$.
\end{proposition}
\begin{proof}
    $\beta\setminus\alpha$ is inhabited
        by \obv{difference_with_proper_subset_is_inhabited}.
    Take $\gamma$ such that $\gamma$ is an \in-minimal element of $\beta\setminus\alpha$
        by \obv{regularity}.
    Now $\gamma\in\beta$ by \obv{setminus_elim_left}.
    Thus $\gamma$ is an ordinal by \obv{elem_of_ordinal_is_ordinal}.
    Hence $\gamma\subseteq\beta$ by \obv{ordinal_elem_implies_subseteq}.
    For all $\delta\in\beta\setminus\alpha$ we have $\delta\notin\gamma$
        by \auto{disjoint}.
    Thus $\gamma\setminus\alpha = \emptyset$
        by \auto{setext,setminus_intro,setminus_elim_left,setminus_elim_right,elem_subseteq}.
    Hence $\gamma\subseteq\alpha$
        by \obv{setminus_eq_emptyset_iff_subseteq}.
    For all $\delta\in\alpha$ we have $\delta\in\gamma$
        by \auto{ordinal_elem_connex,ordinal_is_transitiveset,elem_of_ordinal_is_ordinal,transitiveset_iff_subseteq,elem_subseteq,setminus_elim_right}.
    Thus $\alpha\subseteq\gamma$ by \obv{subseteq}.
    Thus $\alpha = \gamma$ by \obv{subseteq_antisymmetric}.
\end{proof}

\begin{proposition}\label{ordinal_elem_implies_proper_subset}
    Let $\alpha,\beta$ be ordinals.
    Suppose $\alpha\in\beta$.
    Then $\alpha\subset\beta$.
\end{proposition}
\begin{proof}
    Follows by \obv{subset,ordinal_elem_implies_subseteq,in_irrefl}.
\end{proof}

\begin{proposition}\label{ordinal_preceq_implies_subseteq}
Let $\alpha,\beta$ be ordinals.
Suppose $\alpha\precedeseq\beta$.
Then $\alpha\subseteq\beta$.
\end{proposition}
\begin{proof}
\end{proof}

\begin{proposition}\label{ordinal_elem_or_subseteq}
    Let $\alpha,\beta$ be ordinals.
    Then $\alpha\in\beta$ or $\beta\subseteq\alpha$.
\end{proposition}
\begin{proof}
    Follows by \auto{ordinal_elem_connex,ordinal_elem_implies_subseteq,subseteq}.
\end{proof}


\begin{proposition}\label{ordinal_subseteq_or_subseteq}
    Let $\alpha,\beta$ be ordinals.
    Then $\alpha\subseteq\beta$ or $\beta\subseteq\alpha$.
\end{proposition}
\begin{proof}
    Follows by \auto{ordinal_elem_connex,ordinal_elem_implies_subseteq,subseteq}.
\end{proof}


\begin{proposition}\label{ordinal_subseteq_implies_elem_or_eq}
    Let $\alpha,\beta$ be ordinals.
    Suppose $\alpha\subseteq\beta$.
    Then $\alpha\in\beta$ or $\alpha = \beta$.
\end{proposition}
\begin{proof}
    Follows by \auto{subset,ordinal_proper_subset_implies_elem}.
\end{proof}


\begin{corollary}\label{ordinal_subset_trichotomy}
    Let $\alpha,\beta$ be ordinals.
    Then $(\alpha\subset\beta \lor \beta\subset\alpha) \lor \alpha = \beta$.
\end{corollary}
\begin{proof}
    Follows by \auto{subset,ordinal_elem_connex,ordinal_elem_implies_proper_subset}.
\end{proof}

\begin{proposition}\label{ordinal_nor_elem_implies_eq}
    Let $\alpha,\beta$ be ordinals.
    Suppose neither $\alpha\in\beta$ nor $\beta\in\alpha$.
    Then $\alpha = \beta$.
\end{proposition}
\begin{proof}
    Follows by \obv{ordinal_elem_connex}.
\end{proof}

\begin{proposition}\label{ordinal_in_trichotomy}
    Let $\alpha,\beta$ be ordinals. Then $(\alpha\in\beta \lor \beta\in\alpha) \lor \alpha = \beta$.
\end{proposition}
\begin{proof}
    Follows by \obv{ordinal_elem_connex}.
\end{proof}

\begin{corollary}\label{ordinal_prec_trichotomy}
    Let $\alpha,\beta$ be ordinals.
    Suppose neither $\alpha\precedes\beta$ nor $\beta\precedes\alpha$.
    Then $\alpha = \beta$.
\end{corollary}
\begin{proof}
    Follows by \obv{ordinal_nor_elem_implies_eq}.
\end{proof}

\begin{corollary}\label{ordinal_elem_or_superset}
    Let $\alpha,\beta$ be ordinals. Then $\alpha\in \beta$ or $\beta\subseteq \alpha$.
\end{corollary}
\begin{proof}
    Follows by \obv{ordinal_elem_or_subseteq}.
\end{proof}



\subsection{Construction of ordinals}

\begin{proposition}\label{emptyset_is_ordinal}
    $\emptyset$ is an ordinal.
\end{proposition}
\begin{proof}
    Follows by \auto{ordinal,transitiveset,emptyset}.
\end{proof}

% The proof of this theorem benefits from the alternate definition of
% transitivity in terms of $\subseteq$.
\begin{classification}\label{suc_ordinal}
    Let $\alpha$ be an ordinal.
    $\suc{\alpha}$ is an ordinal.
\end{classification}
\begin{proof}
    $\suc{\alpha}$ is \in-transitive by \obv{ordinal,suc_of_transitiveset_is_transitiveset}.
    For every $\beta\in\alpha$ we have that $\beta$ is \in-transitive
        by \obv{ordinal,transitiveset}.
    Follows by \auto{ordinal,suc_iff}.
\end{proof}

\begin{proposition}\label{ordinal_iff_suc_ordinal}
    $\alpha$ is an ordinal iff $\suc{\alpha}$ is an ordinal.
\end{proposition}
\begin{proof}
    Follows by \auto{suc_ordinal,suc_intro_self,elem_of_ordinal_is_ordinal}.
\end{proof}

\begin{proposition}\label{ordinal_in_suc}
    Let $\alpha$ be an ordinal.
    Then $\alpha\in\suc{\alpha}$.
\end{proposition}
\begin{proof}
    Follows by \obv{suc_intro_self}.
\end{proof}

\begin{corollary}\label{ordinal_precedes_suc}
    Let $\alpha$ be an ordinal.
    Then $\alpha\precedes \suc{\alpha}$.
\end{corollary}
\begin{proof}
    Follows by \obv{suc_ordinal,ordinal_in_suc}.
\end{proof}

\begin{proposition}\label{ordinal_elem_implies_subset_of_suc}
    Let $\alpha,\beta$ be ordinals.
    Suppose $\alpha\in\beta$.
    Then $\alpha\subseteq\suc{\beta}$.
\end{proposition}
\begin{proof}
    Follows by \auto{ordinal_elem_implies_subseteq,subseteq_self_suc_intro,subseteq_transitive}.
\end{proof}

\begin{proposition}\label{unions_of_ordinal_is_ordinal}
    Let $\alpha$ be an ordinal.
    Then $\unions{\alpha}$ is an ordinal.
\end{proposition}
\begin{proof}
    For all $x, y$ such that $x\in y\in \unions{\alpha}$ we have $x\in \unions{\alpha}$
        by \auto{unions_intro,unions_iff,transitiveset,ordinal}.
    Thus $\unions{\alpha}$ is \in-transitive by \obv{transitiveset}.
    Every element of $\unions{\alpha}$ is \in-transitive
        by \auto{unions_iff,ordinal,transitiveset}.
    Follows by \obv{ordinal}.
\end{proof}

\begin{proposition}\label{union_of_ordinals_is_ordinal}
    Suppose every element of $A$ is an ordinal.
    Then $\unions{A}$ is an ordinal.
\end{proposition}
\begin{proof}
    Every element of $A$ is an \in-transitive set by \auto{ordinal}.
    Hence $\unions{A}$ is an \in-transitive set
        by \obv{unions_family_of_transitiveset_is_transitiveset}.
    Every element of $\unions{A}$ is an \in-transitive set
        by \auto{unions_iff,ordinal}.
    Follows by \obv{ordinal_intro}.
\end{proof}

\begin{lemma}\label{ordinal_subseteq_unions}
    Let $\alpha$ be an ordinal.
    Then $\unions{\alpha}\subseteq \alpha$.
\end{lemma}
\begin{proof}
    Follows by \obv{ordinal,transitiveset_iff_unions_subseteq}.
\end{proof}

\begin{proposition}\label{union_of_two_ordinals_is_ordinal}
    Let $\alpha,\beta$ be ordinals.
    Then $\alpha\union\beta$ is an ordinal.
\end{proposition}
\begin{proof}
    $\alpha\union\beta$ is \in-transitive by \obv{union_of_transitiveset_is_transitiveset,ordinal}.
    Every element of $\alpha\union\beta$ is \in-transitive
        by \auto{transitiveset,union_iff,ordinal}.
    Follows by \obv{ordinal}.
\end{proof}

\begin{lemma}\label{ordinal_union_below}
    Let $\kappa$ be an ordinal.
    Let $\alpha,\beta\in\kappa$.
    Then $\alpha\union\beta\in\kappa$.
\end{lemma}
\begin{proof}
    $\alpha$ and $\beta$ are ordinals by
        \obv{elem_of_ordinal_is_ordinal}.
    $\alpha\subseteq\beta$ or $\beta\subseteq\alpha$ by
        \obv{ordinal_subseteq_or_subseteq}.
    Hence $\alpha\union\beta=\alpha$ or
        $\alpha\union\beta=\beta$ by
        \obv{union_absorb_subseteq_left,union_absorb_subseteq_right}.
\end{proof}

\begin{lemma}\label{three_ordinal_union_below}
    Let $\kappa$ be an ordinal.
    Let $\alpha,\beta,\gamma\in\kappa$.
    Then $(\alpha\union\beta)\union\gamma\in\kappa$.
\end{lemma}
\begin{proof}
    $\alpha\union\beta\in\kappa$ by \obv{ordinal_union_below}.
    Follows by \obv{ordinal_union_below}.
\end{proof}

\begin{proposition}\label{ordinal_empty_or_emptyset_elem}
    Let $\alpha$ be an ordinal.
    Then $\alpha=\emptyset$ or $\emptyset\in\alpha$.
\end{proposition}
\begin{proof}[Proof by \in-induction on $\alpha$]
    Follows by \auto{setext,emptyset,elem_of_ordinal_is_ordinal,ordinal_transitivity}.
\end{proof}

\begin{proposition}\label{transitive_set_of_ordinals_is_ordinal}
    Let $A$ be a set.
    Suppose that for every $\alpha\in A$ we have $\alpha$ is an ordinal.
    Suppose that $A$ is \in-transitive.
    Then $A$ is an ordinal.
\end{proposition}
\begin{proof}
    Follows by \auto{ordinal}.
\end{proof}

% Apparently Russell first noticed this antinomy while reading a paper by Burali-Forti.
% Typographic NB: Cesare Burali-Forti is a single person, therefore only a single hyphen!
\begin{theorem}[Burali-Forti antinomy]\label{buraliforti_antinomy}
    There exists no set $\Omega$ such that
    for all $\alpha$ we have $\alpha\in \Omega$ iff $\alpha$ is an ordinal.
\end{theorem}
\begin{proof}
    Assume there exists a set $\Omega$ such that
        for all $\alpha$ we have $\alpha\in \Omega$ iff $\alpha$ is an ordinal.
    Take $\Omega$ such that for all $\alpha$ we have
        $\alpha\in \Omega$ iff $\alpha$ is an ordinal.
    For all $x, y$ such that $x\in y\in \Omega$ we have $x\in \Omega$
        by \auto{elem_of_ordinal_is_ordinal}.
    Thus $\Omega$ is \in-transitive by \obv{transitiveset}.
    For every $\alpha\in\Omega$ we have $\alpha$ is an ordinal.
    Thus $\Omega$ is an ordinal by \auto{ordinal}.
    Therefore $\Omega\in\Omega$.
    Contradiction by \obv{in_irrefl}.
\end{proof}

\begin{proposition}\label{inters_of_ordinals_is_ordinal}
    Let $A$ be an inhabited set.
    Suppose for every $\alpha\in A$ we have $\alpha$ is an ordinal.
    Then $\inters{A}$ is an ordinal.
\end{proposition}
\begin{proof}
    For all $x,y$ such that $x\in y\in\inters{A}$ we have $x\in\inters{A}$
        by \auto{inters_iff_forall,ordinal_transitivity}.
    Thus $\inters{A}$ is \in-transitive by \obv{transitiveset}.
    Take $\alpha$ such that $\alpha\in A$.
    For every $y\in\inters{A}$ we have $y$ is an ordinal
        by \obv{inters_iff_forall,elem_of_ordinal_is_ordinal}.
    Follows by \auto{ordinal}.
\end{proof}

\begin{proposition}\label{inters_of_ordinals_subseteq}
    Let $A$ be an inhabited set.
    Suppose for every $\alpha\in A$ we have $\alpha$ is an ordinal.
    Then for all $\alpha\in A$ we have $\inters{A}\subseteq \alpha$.
\end{proposition}
\begin{proof}
    Follows by \obv{inters_subseteq_elem}.
\end{proof}

\begin{proposition}\label{inters_of_ordinals_elem}
    Let $A$ be an inhabited set.
    Suppose for every $\alpha\in A$ we have $\alpha$ is an ordinal.
    Then $\inters{A}\in A$.
\end{proposition}
\begin{proof}
    Follows by \auto{inters_of_ordinals_is_ordinal,in_irrefl,inters_iff_forall,ordinal_subseteq_implies_elem_or_eq,inters_subseteq_elem}.
\end{proof}

\begin{proposition}\label{inters_of_ordinals_is_minimal}
    Let $A$ be an inhabited set.
    Suppose for every $\alpha\in A$ we have $\alpha$ is an ordinal.
    Then $\inters{A}$ is an \in-minimal element of $A$.
\end{proposition}
\begin{proof}
    $\inters{A}\in A$ by \obv{inters_of_ordinals_elem}.
    For all $\alpha\in A$ we have $\inters{A}\subseteq \alpha$
        by \obv{inters_of_ordinals_subseteq}.
    For all $\delta\in\inters{A}$ we have $\delta\notin A$
        by \auto{elem_of_ordinal_is_ordinal,ordinal_elem_implies_proper_subset,subset,subseteq_antisymmetric}.
    Follows by \auto{disjoint}.
\end{proof}

\begin{proposition}\label{inters_of_ordinals_is_minimal_alternate}
    Let $A$ be an inhabited set.
    Suppose for every $\alpha\in A$ we have $\alpha$ is an ordinal.
    Then for all $\alpha\in A$ we have $\inters{A} = \alpha$ or $\inters{A}\in\alpha$.
\end{proposition}
\begin{proof}
    Follows by \auto{inters_of_ordinals_is_ordinal,inters_of_ordinals_subseteq,ordinal_subseteq_implies_elem_or_eq}.
\end{proof}

\begin{proposition}\label{inter_of_two_ordinals_is_ordinal}
    Let $\alpha,\beta$ be ordinals.
    Then $\alpha\inter\beta$ is an ordinal.
    \end{proposition}
\begin{proof}
    $\alpha\inter\beta$ is \in-transitive by \auto{inter,ordinal,transitiveset}.
    Every element of $\alpha\inter\beta$ is \in-transitive
        by \obv{transitiveset,inter,ordinal}.
    Follows by \obv{ordinal}.
\end{proof}


\subsection{Limit and successor ordinals}


\begin{definition}\label{limit_ordinal}
    $\lambda$ is a limit ordinal iff
    $\emptyset\precedes \lambda$ % by definition of \precedes this means that $\lambda$ is an ordinal
    and for all $\alpha\in \lambda$ we have $\suc{\alpha}\in \lambda$.
\end{definition}

\begin{definition}\label{successor_ordinal}
    $\alpha$ is a successor ordinal iff
    there exists an ordinal $\beta$ such that $\alpha = \suc{\beta}$.
\end{definition}

\begin{lemma}\label{ordinal_successor_data}
    Let $\beta$ be an ordinal.
    Then $\suc{\beta}$ is a successor ordinal and
        $\unions{\suc{\beta}}=\beta$.
\end{lemma}
\begin{proof}
    $\suc{\beta}$ is a successor ordinal by
        \obv{successor_ordinal}.
    $\beta$ is \in-transitive by
        \obv{ordinal_is_transitiveset}.
    Hence $\unions{\suc{\beta}}=\beta$ by
        \obv{transitiveset_iff_unions_suc}.
\end{proof}

\begin{lemma}\label{successor_ordinal_eq_suc_unions}
    Let $\alpha$ be a successor ordinal.
    Then $\alpha=\suc{\unions{\alpha}}$.
\end{lemma}
\begin{proof}
    Take an ordinal $\beta$ such that $\alpha=\suc{\beta}$ by
        \obv{successor_ordinal}.
    $\beta$ is \in-transitive by
        \obv{ordinal_is_transitiveset}.
    Hence $\unions{\alpha}=\beta$ by
        \auto{transitiveset_iff_unions_suc}.
\end{proof}

\begin{lemma}\label{positive_ordinal_is_limit_or_successor}
    Let $\alpha$ be an ordinal such that $\emptyset\precedes\alpha$.
    Then $\alpha$ is a limit ordinal or $\alpha$ is a successor ordinal.
\end{lemma}
\begin{proof}
    Assume $\alpha$ is not a limit ordinal.
    Take $\beta$ such that $\beta\in\alpha$ and $\suc{\beta}\notin\alpha$
        by \auto{limit_ordinal}.
    Thus $\beta$ is an ordinal by \obv{elem_of_ordinal_is_ordinal}.
    Thus $\suc{\beta}$ is an ordinal by \obv{suc_ordinal}.
    We have $\alpha=\suc{\beta}$
        by \auto{ordinal_elem_connex,suc_elim,in_asymmetric,in_irrefl}.
    Follows by \obv{successor_ordinal}.
\end{proof}

\begin{lemma}\label{zero_not_successorordinal}
    $\emptyset$ is not a successor ordinal.
\end{lemma}
\begin{proof}
    Follows by \auto{successor_ordinal,suc_neq_emptyset}.
\end{proof}

\begin{lemma}\label{zero_not_limitordinal}
    $\emptyset$ is not a limit ordinal.
\end{lemma}
\begin{proof}
    Follows by \obv{limit_ordinal,emptyset,in_irrefl}.
\end{proof}

\begin{lemma}\label{successor_ordinal_is_not_limit_ordinal}
    Let $\alpha$ be an ordinal.
    Then $\suc{\alpha}$ is not a limit ordinal.
\end{lemma}
\begin{proof}
    Suppose $\suc{\alpha}$ is a limit ordinal.
    $\alpha\in\suc{\alpha}$ by \obv{suc_intro_self}.
    Hence $\suc{\alpha}\in\suc{\alpha}$ by \obv{limit_ordinal}.
    Contradiction by \obv{in_irrefl}.
\end{proof}

\begin{lemma}\label{limit_ordinal_is_not_successor_ordinal}
    Let $\alpha$ be a limit ordinal.
    Then $\alpha$ is not a successor ordinal.
\end{lemma}
\begin{proof}
    Suppose $\alpha$ is a successor ordinal.
    Take an ordinal $\beta$ such that $\alpha=\suc{\beta}$ by
        \obv{successor_ordinal}.
    $\suc{\beta}$ is not a limit ordinal by
        \obv{successor_ordinal_is_not_limit_ordinal}.
    Hence $\alpha$ is not a limit ordinal.
    Contradiction.
\end{proof}

\begin{lemma}\label{suc_elem_limitordinal}
    Let $\lambda$ be a limit ordinal.
    Let $\alpha\in\lambda$.
    Then $\suc{\alpha}\in\lambda$.
\end{lemma}
\begin{proof}
    Follows by \obv{limit_ordinal}.
\end{proof}

\begin{lemma}\label{limitordinal_eq_unions}
    Let $\lambda$ be a limit ordinal.
    Then $\unions{\lambda} = \lambda$.
\end{lemma}
\begin{proof}
    $\unions{\lambda}\subseteq\lambda$ by \obv{limit_ordinal,ordinal_subseteq_unions}.
    For all $\alpha\in\lambda$ we have $\alpha\in\suc{\alpha}\in\lambda$ by \obv{suc_intro_self,suc_elem_limitordinal}.
    Thus $\lambda\subseteq\unions{\lambda}$ by \auto{subseteq,unions_intro}.
    Follows by \obv{subseteq_antisymmetric}.
\end{proof}
